\begin{derivation}[从\eqnrefs{eq:sm-one-loop-determinant-master-dp68}到\eqnrefs{eq:sm-coleman-weinberg-si-dp12}]
正文\eqnrefs{eq:sm-one-loop-determinant-master-dp68} 已把常数 Higgs 背景上的二次涨落统一为函数行列式。先取一个实玻色自由度，记其背景质量能量平方为 $\mathcal M^2$：
\begin{equation*}
 V_1^{(B)}(\mathcal M^2)
 =\frac{\hbar c}{2}
 \left(\frac{\mu_{\rm DR}}{\hbar}\right)^{2\epsilon}
 \int\frac{\dd^dk}{(2\pi)^d}
 \ln\!\left[k^2+\frac{\mathcal M^2}{(\hbar c)^2}\right],
 \qquad d=4-2\epsilon.
\end{equation*}
直接积分对数不方便；先对 $\mathcal M^2$ 求导：
\begin{equation*}
 \frac{\partial V_1^{(B)}}{\partial\mathcal M^2}
 =\frac{1}{2\hbar c}
 \left(\frac{\mu_{\rm DR}}{\hbar}\right)^{2\epsilon}
 \int\frac{\dd^dk}{(2\pi)^d}
 \frac{1}{k^2+\mathcal M^2/(\hbar c)^2}.
\end{equation*}
欧氏母积分
\begin{equation*}
 \int\frac{\dd^dk}{(2\pi)^d}\frac1{k^2+m^2}
 =\frac{\Gamma(1-d/2)}{(4\pi)^{d/2}}(m^2)^{d/2-1}
\end{equation*}
给出
\begin{equation*}
 \frac{\partial V_1^{(B)}}{\partial\mathcal M^2}
 =\frac{\mathcal M^2}{2(4\pi)^2(\hbar c)^3}
 \Gamma(-1+\epsilon)(4\pi)^\epsilon
 \left(\frac{E_R^2}{\mathcal M^2}\right)^\epsilon,
 \qquad E_R=c\mu_{\rm DR}.
\end{equation*}
把 $\mathcal M^2$ 从零积分到当前背景值；与背景无关的积分常数只平移真空能，可吸收到宇宙学常数型反项。利用
\begin{align*}
 \Gamma(-2+\epsilon)
 &=\frac1{2\epsilon}+\frac34-\frac{\gamma_E}{2}+\order(\epsilon),\\
 (4\pi)^\epsilon
 &=1+\epsilon\ln4\pi+\order(\epsilon^2),\\
 \left(\frac{E_R^2}{\mathcal M^2}\right)^\epsilon
 &=1+\epsilon\ln\frac{E_R^2}{\mathcal M^2}+\order(\epsilon^2),
\end{align*}
可把未减除结果整理为
\begin{equation*}
 V_1^{(B)}
 =\frac{\mathcal M^4}{64\pi^2(\hbar c)^3}
 \left[
 -\frac1{\bar\epsilon}
 +\ln\frac{\mathcal M^2}{E_R^2}-\frac32
 \right]+\order(\epsilon),
 \qquad
 \frac1{\bar\epsilon}\equiv\frac1\epsilon-\gamma_E+\ln4\pi.
\end{equation*}
$\overline{\rm MS}$ 反项只减去 $1/\bar\epsilon$，所以每个实标量自由度留下
\begin{equation*}
 V_{1,\rm ren}^{(B)}
 =\frac{\mathcal M^4}{64\pi^2(\hbar c)^3}
 \left(\ln\frac{\mathcal M^2}{E_R^2}-\frac32\right).
\end{equation*}

费米高斯积分产生 $\det\mathcal D_F$ 而不是 $\det^{-1/2}\mathcal D_B$，因此闭合费米子环带负号；Dirac 自旋、颜色等多重度再乘入 $n_i$。质量矢量场的常数 $5/6$ 还不能简单从标量结果“照抄”。在维数正则化中，横向 massive-vector 行列式的极化数是 $d-1=3-2\epsilon$，因此其裸贡献含
\begin{align*}
 (3-2\epsilon)
 \left[-\frac1{\bar\epsilon}
 +\ln\frac{\mathcal M_V^2}{E_R^2}-\frac32\right]
 &=-\frac3{\bar\epsilon}
 +3\ln\frac{\mathcal M_V^2}{E_R^2}
 -\frac92+2+\order(\epsilon)\\
 &=-\frac3{\bar\epsilon}
 +3\left[
 \ln\frac{\mathcal M_V^2}{E_R^2}-\frac56
 \right]+\order(\epsilon).
\end{align*}
这里有限的 $+2$ 正是 $(-2\epsilon)$ 乘紫外极点产生的项；它把每个物理矢量自由度的常数从 $3/2$ 改为 $5/6$。在正文采用的 Landau 规范组织下，Goldstone 与 ghost 的背景依赖另外计入各自的 $n_i,\mathcal M_i$，不会再重复计入这一横向矢量多重度。

把所有粒子种类统一写成带符号的 $n_i$，并加回树级势，最终得到
\begin{equation*}
 V_{\rm eff}=V_0+\frac1{64\pi^2(\hbar c)^3}
 \sum_i n_i\mathcal M_i^4
 \left[\ln\frac{\mathcal M_i^2}{E_R^2}-c_i\right]
 +\order(\hbar^2),
\end{equation*}
其中标量和费米子 $c_i=3/2$，横向质量矢量 $c_i=5/6$，即正文\eqnrefs{eq:sm-coleman-weinberg-si-dp12}。每项的 $\mathcal M_i^4/(\hbar c)^3$ 均具有 $\mathrm{J\,m^{-3}}$ 量纲。
\end{derivation}
